\documentclass[11pt]{article}
% Shared preamble for all daily exercises
% Update this file to change settings (padding, parindent, etc.) for all exercises

\usepackage[margin=0.75in]{geometry}
\usepackage{amsmath}
\usepackage{amssymb}

% Custom settings
\setlength{\parindent}{0pt}
\setlength{\parskip}{0.2cm}

\title{Daily Exercise}
\author{Dhruman Gupta}
\date{18th April}

\begin{document}

% Common header for daily exercises
% This file can be included in other LaTeX documents

\maketitle

\noindent
\textbf{Course:} CS-3330, Theory of Computation

\vspace{0.2cm}

\noindent
\textbf{Question:}\noindent 
An oracle Turing machine is a Turing machine with an extra tape called the oracle tape and three special states: QUERY, YES, NO. When the machine enters QUERY, it receives an answer in one step: YES if the string on the oracle tape belongs to the oracle set $A$, else NO. For a complexity class $C$ and an oracle set $A$, $C^A$ denotes the class of languages decidable by a machine of type $C$ with access to oracle $A$.

Show that $P^{SAT}$ contains both $NP$ and $coNP$.

\textbf{Answer:}

We will show that both $NP \subseteq P^{SAT}$ and $coNP \subseteq P^{SAT}$.

First, let $L \in NP$. Then, by Cook-Levin, there exists a polynomial time reduction from $L$ to $SAT$. That is, for every input $x$, we can construct in polynomial time a Boolean formula $\phi_x$ such that
\[
x \in L \iff \phi_x \in SAT.
\]

Now construct a deterministic polynomial time oracle Turing machine with oracle $SAT$. On input $x$, it does the following:
\begin{itemize}
    \item Construct $\phi_x$ in polynomial time.
    \item Write $\phi_x$ on the oracle tape and enter the QUERY state.
    \item If the oracle answers YES, accept. Otherwise, reject.
\end{itemize}

This machine runs in polynomial time, since constructing $\phi_x$ takes polynomial time and the oracle query takes one step. Hence $L \in P^{SAT}$. Therefore,
\[
NP \subseteq P^{SAT}.
\]

Now let $L \in coNP$. Then $\overline{L} \in NP$. By the above argument,
\[
\overline{L} \in P^{SAT}.
\]

So there exists a deterministic polynomial time oracle machine with oracle $SAT$ deciding $\overline{L}$. To decide $L$, just run that machine and reverse the final answer: accept if it rejects, and reject if it accepts. This is still a deterministic polynomial time oracle machine with oracle $SAT$.

Hence $L \in P^{SAT}$, and so
\[
coNP \subseteq P^{SAT}.
\]

Therefore, $P^{SAT}$ contains both $NP$ and $coNP$.

\end{document}
